\documentclass[10pt,a4paper]{amsart}
\usepackage[margin=19mm]{geometry}
\usepackage{amsmath,amssymb,mathtools}
\usepackage[scr=rsfs,cal=euler]{mathalfa}
\usepackage{xfrac}
\usepackage[unicode,hidelinks,bookmarks=false]{hyperref}
\newcommand{\ie}{i.e.\ }
\newcommand{\cf}{cf.\ }
\newcommand{\resp}{respectively\ }
\newcommand{\Subsection}{\subsection}
\providecommand{\nobreakdash}{\nobreak\hbox{-}}
\setlength{\emergencystretch}{2em}
\multlinegap0pt
\usepackage{amssymb}
\usepackage{mathtools}
\usepackage{xcolor}
\usepackage[shortlabels]{enumitem}
\newtheorem{proposition}{Proposition}
\usepackage{cleveref}
\crefname{proposition}{Proposition}{Propositions}
\newcommand{\I}{\mathscr{I}}
\newcommand{\K}{\mathscr{K}}
\newcommand{\gbd}{\propto_\Gamma}
\title{On AF- and type I-ideals in certain crossed product C*-algebras}
\author{Alexander G. Ravnanger}
\date{Source: arXiv:2609.11297v1 (CC BY 4.0)}
\begin{document}
\maketitle

\noindent\emph{Source and licence.} On AF- and type I-ideals in certain crossed product C*-algebras. Version arXiv:2609.11297v1 (CC BY 4.0), distributed by arXiv under the Creative Commons Attribution 4.0 International licence, http://creativecommons.org/licenses/by/4.0/, as recorded in the arXiv licence metadata for this exact version and shown on the abstract page https://arxiv.org/abs/2609.11297v1. Full licence notice: \texttt{licenses/arxiv-2609.11297-CC-BY-4.0.txt}.\par\smallskip

\noindent\textit{Editorial scope.} Counted unit: the article's proposition proposition (source0.tex lines 468--470). The article's complete original proof of it is delivered with it (source0.tex lines 472--490, one \texttt{proof} environment of 19 lines). No mathematics is written, repaired or re-derived: the statement and the proof are the article's own lines, in the article's own notation, and no retained line is substituted or re-flowed.  Retained uncounted as the source-local context the proof needs: lines 393--402.  Nothing was omitted from any retained span, so the package carries no omission record.  No retained line contains a citation, so the wrapper carries no bibliography and no bibliography entry.  No retained line contains a figure environment, an image-inclusion command or drawing code, so no image is delivered and the package declares no asset dependency.  Editorial formatting only, in addition: an AMS \texttt{amsart} preamble that restates the article's own declarations and macros used by the retained lines, namely the environments \texttt{proposition} on separate counters, together with the standard packages those lines need.  The retained environments print in this document as Proposition 1 in order of appearance, whereas the article prints the same items with the numbering of its own full document.  The complete native source of the article as distributed in the arXiv e-print for this exact version is bundled unchanged as \texttt{sources/210003/source0.tex}.
	\begin{proposition}
		\label{prop_compact}
		Let \(\Gamma\) denote an exact group. The compact ideals in \(C_u^*(\Gamma)\) are precisely the ideals of the form \(\I_A\) for subsets \(A \subseteq \Gamma\). Moreover, every ideal in \(C_u^*(\Gamma)\) is an increasing union of such compact ideals. 
		
		For two subsets \(A, B \subseteq \Gamma\), we have the following: 
		\begin{enumerate}
			\item \(\I_A \subseteq \I_B\) if and only if \(A \gbd B\). 
			\item \(\I_{A \cup B} = \I_A + \I_B\). 
		\end{enumerate}
	\end{proposition}

	\begin{proposition}
		If \(A \subseteq \Gamma\) is a finite union of sparse subsets of a countable group, then \(\I_A/ \K\) is type \(I_0\). In particular, \(\I_A\) is type I. 
	\end{proposition}

	\begin{proof} 
		Clearly the last statement follows from the first as an extension of two type I C*-algebras is again type I. Let \(\pi \colon \I_A \to \I_A/\K\) denote the quotient map. It suffices to assume that \(A\) is sparse and show that for any pair of elements \(g,h \in \Gamma\), finite subset \(F \subseteq \Gamma\) and corresponding \(f_k \in C_0(X_A)\) for \(k \in F\), we have 
		\begin{equation}
			\label{eq_sparse}
		\pi \left( (1_{hA}u_g)^* \left( \sum_{k \in F} f_k u_k \right) (1_{hA}u_g) \right) \in \pi(C_0(X_A))
		\end{equation}
		Indeed, that would show that the element \(1_{hA}u_g\) is abelian modulo \(\K\). Moreover, if \(C\) is any compact-open set contained in \(hA\), the hereditary C*-subalgebra generated by \(1_Cu_g\) is contained in the one generated by \(1_A u_g\), and so we conclude that for all such \(C\), the element \(1_Cu_g\) is abelian modulo \(\K\). Letting these run through all compact-open sets of \(X_A\) and all \(g \in \Gamma\), we conclude that \(\I_A\) is generated as a C*-algebra by elements that are abelian modulo \(\K\), assuming that \(A\) is sparse. If \(A\) is a finite union of sparse sets, we combine the previous case with \cref{prop_compact}.
		
		To verify \cref{eq_sparse}, we expand the brackets in the element inside \(\pi\): 
		\begin{align*}
			(1_{hA}u_g)^* \left( \sum_{k \in F} f_k u_k \right) (1_{hA}u_g) &= \sum_{k \in F} u_g^* 1_{hA} f_k u_k 1_{hA} u_g \\
			&= \sum_{k \in F} u_g^* 1_{hA} f_k 1_{khA} u_{kg} \\
			&= \sum_{k \in F} 1_{g^{-1} hA} \alpha_g^{-1}(f_k) 1_{g^{-1}khA} u_{g^{-1}k g} \\
			&= \sum_{k \in F} 1_{g^{-1}h(A \cap h^{-1}khA)} \alpha_g^{-1}(f_k) u_{g^{-1}kg}
		\end{align*}
		By assumption, \(A \cap h^{-1}khA\) is finite unless \(k = e\). Hence, 
		\[	\pi \left( (1_{hA}u_g)^* \left( \sum_{k \in F} f_k u_k \right) (1_{hA}u_g) \right) = \pi(1_{g^{-1}hA} \alpha_g^{-1}(f_e) u_e) \in \pi(C_0(X_A)),\]
		which completes the proof. 
	\end{proof}

\end{document}
